\section{Overview of TAC Industries}

TAC Industries is a security-centric investigative infrastructure and case management platform designed to model how sensitive intelligence and investigative workflows can be supported within a controlled enterprise environment. The project does not represent a real-world law enforcement system, but rather a fictionalised and ethically constrained platform inspired by established investigative and intelligence system design principles.

At its core, TAC Industries combines secure infrastructure design with a role-aware case and person registry. The system is intended to demonstrate how modern enterprise security concepts, including Zero Trust networking, identity-centric access control, and audited privileged access can be applied cohesively to an investigative use case. Rather than focusing on feature-rich application functionality, the project prioritises architectural correctness, access governance, and data protection from the outset.


\section{Intended Users and Roles}

TAC Industries is designed to support a small number of clearly defined user roles, each with strictly limited permissions aligned to the principle of least privilege. These roles are deliberately separated to reduce insider risk and prevent any single user from exercising excessive control over both infrastructure and sensitive data.

Operational users are responsible for creating and reading investigative case data within their authorised scope. Intelligence users have read-only access to case summaries and reports, enabling situational awareness without the ability to modify or introduce data. Lead or supervisory users are granted elevated data authority to manage case lifecycle actions such as approval, closure, and correction, with all actions subject to audit logging.

Infrastructure and systems administration is intentionally segregated from investigative data access. Administrative users manage identity services, servers, monitoring, and endpoints, but are technically restricted from accessing case content. This separation of duties ensures that infrastructure control does not imply access to sensitive intelligence or personal information.

\section{Ethical and Legal Context}

Given the sensitive nature of investigative and intelligence-style systems, ethical and legal considerations are treated as a first-class design concern within TAC Industries. The platform operates exclusively on synthetic data and does not process real personal information. All identifiers, documents, and records are artificially generated for demonstration purposes using LLM's that will be referenced throughout the project.

The system is designed to minimise data collection, avoid automated profiling or decision-making, and clearly distinguish between factual records and indicators or flags. Access to sensitive information is strictly role-controlled, and all data modifications are logged to ensure traceability and accountability \cite{gdpr2016}.

All monitoring, inspection, and identity enforcement mechanisms are implemented within a controlled academic environment using synthetic data, with explicit boundaries defined to avoid intrusive monitoring of personal or guest traffic. By embedding these considerations directly into the architecture, the project demonstrates how ethical responsibility and legal awareness can be integrated into secure system design rather than treated as external constraints.